\clearpage\begingroup
\SourceLeaf{148}\NoSourceHead
\makeatletter\tagsleft@true\let\veqno\@@leqno\makeatother
\begin{center}
[VI.]\par\bigskip
\textbf{REFLEXION VON DER HINTERN FLÄCHE EINER LINSE.}\par\medskip
\rule{.44\linewidth}{.4pt}\par
[Aus Handbuch 19, Be, S. 230.]\par
\rule{.44\linewidth}{.4pt}
\end{center}
\begin{center}
\begin{tikzpicture}[x=.011cm,y=.011cm,line width=.45pt]
\draw (48,0)--(1015,0);
\draw (187,172) .. controls (139,128) and (113,57) .. (113,0)
 .. controls (113,-65) and (144,-130) .. (188,-173);
\draw (187,172) .. controls (211,99) and (219,52) .. (219,0)
 .. controls (219,-65) and (210,-127) .. (188,-173);
\draw (142,6)--(142,-6);
\draw (434,10)--(434,-10);
\draw (962,10)--(962,-10);
\node[below,inner sep=2pt] at (142,-5) {$E$};
\node[below,inner sep=2pt] at (235,-5) {$N$};
\node[below,inner sep=2pt] at (434,-5) {$F$};
\node[below,inner sep=2pt] at (434,-44) {\itshape Brennpunkt};
\node[below,inner sep=2pt] at (962,-5) {$E'$};
\end{tikzpicture}
\end{center}
\noindent\begin{tabular}{@{}l@{\quad}p{.91\linewidth}@{}}
$r$&Halbmesser der Vorderfläche, welche die Axe in $N$ trifft,\\
$r^{*}$&Halbmesser der Hinterfläche, an welcher die Reflexion geschieht,
\end{tabular}
\[
\frac{r}{n-1}=f,\quad\frac{r^{*}}{n}=f^{*},
\]
\noindent Dicke $=ne$.

In Beziehung auf das System von zwei Brechungen und Eine Reflexion,
\begin{center}
$F$ Brennpunkt, $E$ erster, $E'$ zweiter Hauptpunkt,\\
$E'-E=2F=$ doppelte Brennweite.
\end{center}
\noindent Die Coordinaten werden als wachsend betrachtet in dem Sinn der Bewegung der rückkehrenden Lichtstrahlen.
\begin{center}\textbf{\textls[120]{Formeln.}}\end{center}
\begin{align}
E&=N-\frac{ef}{f-e}\tag*{I.}\\[5pt]
E'&=N+\frac{(f^{*}-e)f}{f^{*}+f-e}\tag*{II.}\\[5pt]
E'-E&=2F=\frac{fff^{*}}{(f-e)(f+f^{*}-e)}.\tag*{III.}
\end{align}
\begin{center}\rule{.25\linewidth}{.4pt}\end{center}
\clearpage\SourceLeaf{149}\HistoricalHead{144}{L}{PHYSIK.\quad NACHLASS.}
\begin{center}\textbf{ERLÄUTERUNGEN.}\end{center}
\noindent\hspace*{\parindent}Die gestellte Aufgabe ist ein Beispiel zur Theorie der Hauptpunkte eines Linsensystems, das etwas allgemeiner ist als das von \textsc{Gauss} im Art. 11 seiner \textit{Dioptrischen Untersuchungen} (Werke V, S. 261 ff.) behandelte.

Man erhält die \textsc{Gauss}schen Formeln I, II, III, indem man in der Bezeichnung der \textit{Dioptrischen Untersuchungen} setzt:
\[
n^0=1,\quad n'=n,\quad n''=-n,\quad n^{*}=-1.
\]
\noindent Dann ergeben sich die dort (S. 250) definierten Größen $u$ und $t$:
\[
\begin{gathered}
u^0=\frac1f,\quad u'=\frac2{f^{*}},\quad u''=+\frac1f,\\[5pt]
t'=-e,\quad t''=-e,
\end{gathered}
\]
\noindent und daraus folgen dann die \textsc{Gauss}schen Koeffizienten $g,k,l$ (Werke V, S. 250, Gl. 5):
\[
\begin{aligned}
g&=\frac{ff^{*}-2ef-2ef^{*}+2e^2}{ff^{*}},\\[5pt]
k&=-2\frac{(f-e)^2+f^{*}(f-e)}{fff^{*}},\\[5pt]
l&=\frac{2ef^{*}-2e^2-ff^{*}+2ef}{ff^{*}},
\end{aligned}
\]
\noindent aus denen sich weiter gemäß Tabelle I der \textit{Dioptrischen Untersuchungen} (Werke V, S. 255) die angegebenen Werte für $E$ und $E'$ berechnen lassen. Die genannte Tabelle enthält übrigens einen Druckfehler, indem es heißen muß: $E=N^0-\dfrac{n^0(1-e)}k$.
\begin{flushright}\textsc{Schaefer.}\end{flushright}
\begin{center}\rule{.30\linewidth}{.4pt}\end{center}
\clearpage\endgroup
